\documentclass[11pt]{article}

\usepackage[margin=1.7cm]{geometry}
\usepackage{amsmath}
\usepackage{amssymb}
\usepackage{enumitem}
\usepackage{titlesec}
\usepackage{xcolor}
\usepackage{fancyhdr}
\usepackage{array}
\usepackage{makecell}

\titleformat{\section}{\large\bfseries}{\thesection}{0.6em}{}
\titlespacing{\section}{0pt}{1.1em}{0.5em}

\setlength{\headheight}{14pt}

\pagestyle{fancy}
\fancyhf{}
\fancyhead[L]{Analysing Two-Way Tables}
\fancyfoot[C]{\thepage}
\renewcommand{\headrulewidth}{0.4pt}

\fancypagestyle{firstpage}{%
  \fancyhf{}
  \fancyhead[L]{Analysing Two-Way Tables}
  \fancyhead[R]{Name: \makebox[4cm]{\hrulefill}}
  \fancyfoot[C]{\thepage}
  \renewcommand{\headrulewidth}{0.4pt}
}

\newcommand{\ans}{\hrulefill}
\newcommand{\anshere}{\par\vspace{0.3em}\noindent\ans\par\vspace{0.3em}}
\setlist[enumerate,1]{leftmargin=1.4em, itemsep=0.5em, topsep=0.2em}
\setlist[enumerate,2]{leftmargin=1.3em, itemsep=0.3em, topsep=0.2em}

\newcolumntype{Y}{>{\centering\arraybackslash}p{1.9cm}}
\newcolumntype{L}{>{\raggedright\arraybackslash}p{3.3cm}}

\begin{document}
\thispagestyle{firstpage}

\begin{center}
  {\Large\bfseries Analysing Two-Way Tables}\\[2pt]
  {\normalsize NCEA Level 2 Statistics --- Bivariate Categorical Data}
\end{center}

\vspace{0.4em}
\noindent A \textbf{two-way table} sorts data by two categories at once. Whenever the
\textbf{groups being compared are different sizes}, comparing raw counts can be
misleading --- you need to convert counts to \textbf{percentages} first.

\vspace{0.3em}
\noindent\textbf{Row percentage} = $\dfrac{\text{cell value}}{\text{row total}} \times 100$.
This tells you, \emph{within one group}, what share falls into each category --- which
is exactly what you need to fairly compare groups of different sizes.

\vspace{0.4em}
\noindent\textbf{Worked example:} 40 students at School A and 200 students at School B
were asked if they play a sport. 28 students at School A said yes, and 120 students at
School B said yes. More School B students said yes (120 vs 28) --- but School B also
asked far more students, so this comparison is misleading. Using row percentages instead:
\[\text{School A: } \frac{28}{40}\times 100 = 70\%, \qquad
  \text{School B: } \frac{120}{200}\times 100 = 60\%.\]
A \emph{higher percentage} of School A students play a sport, even though fewer of them
said yes in total.

\section{Calculating Percentages}
\noindent 170 Year 13 students were asked how they studied for an exam, and whether they
passed.

\begin{center}
\begin{tabular}{|L|Y|Y|Y|}
\hline
 & \textbf{Passed} & \textbf{Failed} & \textbf{Total} \\\hline
\textbf{Group study}    & 56 & 14 & \\\hline
\textbf{Solo study}     & 42 & 18 & \\\hline
\textbf{Did not study}  & 10 & 30 & \\\hline
\textbf{Total}          &    &    & \\\hline
\end{tabular}
\end{center}

\begin{enumerate}
  \item Copy and complete the table, filling in all the totals.
  \item What percentage of students who studied in a group passed?
  \anshere
  \item What percentage of students who studied alone passed?
  \anshere
  \item What percentage of students who did not study passed?
  \anshere
  \item What percentage of \emph{all} 170 students passed?
  \anshere
\end{enumerate}

\section{Comparing Distributions}
\noindent Use your percentages from Section 1 to answer these.

\begin{enumerate}
  \item Write a sentence comparing the pass rate of students who studied in a group with
    students who did not study, using percentages.
  \anshere
  \item Tyler says, ``Group study must be better than solo study, because more students
    who did group study passed (56) than students who did solo study (18 failed vs 14
    failed).'' Explain what is wrong with the way Tyler is comparing these two groups.
  \anshere
  \item Using percentages, write a correct comparison between group study and solo
    study.
  \anshere
\end{enumerate}

\section{Describing a Trend}
\noindent 400 people, in four age groups, were asked which type of transport they use
most often to get to work.

\begin{center}
\begin{tabular}{|L|Y|Y|Y|Y|}
\hline
 & \textbf{Public transport} & \textbf{Car} & \textbf{Other} & \textbf{Total} \\\hline
\textbf{15--24} & 45 & 30 & 25 & \\\hline
\textbf{25--44} & 30 & 60 & 10 & \\\hline
\textbf{45--64} & 20 & 70 & 10 & \\\hline
\textbf{65+}    & 15 & 55 & 30 & \\\hline
\end{tabular}
\end{center}

\begin{enumerate}
  \item Copy and complete the table, filling in the total for each age group.
  \item Calculate the percentage of each age group that uses public transport most
    often.
  \anshere
  \item Describe the trend in public transport use as age increases. Refer to your
    percentages.
  \anshere
  \item The percentage using ``Other'' transport is much higher for the 65+ age group
    than for the 45--64 age group. Suggest a reason this might be the case.
  \anshere
\end{enumerate}

\section{Making a Justified Conclusion}
\noindent Two Year 12 classes were asked whether they like rugby.

\begin{center}
\begin{tabular}{|L|Y|Y|Y|}
\hline
 & \textbf{Like} & \textbf{Dislike} & \textbf{Total} \\\hline
\textbf{Class A} & 14 & 8  & \\\hline
\textbf{Class B} & 11 & 13 & \\\hline
\end{tabular}
\end{center}

\begin{enumerate}
  \item Calculate the percentage of each class that likes rugby.
  \anshere
  \item Class A's percentage is noticeably higher than Class B's. Explain why this
    difference might \emph{not} show a real difference between the two classes, given
    how many students are in each sample.
  \anshere
  \item What could you do to become more confident that there is a genuine difference
    between the two classes, rather than the difference being due to chance?
  \anshere
  \item In Section 3, you found that public transport use decreases as age increases.
    This survey only sampled people from one city. Explain why you should be cautious
    about concluding this trend is true for \emph{everyone}, everywhere.
  \anshere
\end{enumerate}

\end{document}
